\scanpage{32}
\mspage{6}
Let $y$ be in $Y$. Then $\theta(y)$ is in the norm closed convex hull of $\{t(y):t\in T\}$ by [3, Th. 13, p. 422]. Therefore, given any $\varepsilon>0$ there exist $t_1,\ldots,t_n$ in $T$ and $\alpha_i\geqq0$ $(i=1,\ldots,n)$ such that $\sum\alpha_i=1$ and
\[
\|\theta(y)=\sum\alpha_it_i(y)\|<\varepsilon/(M+1),
\]
where $M=\sup\{\|t\|:t\in T\}$ (which is finite by the uniform boundedness principle.) Then if $t\geqq t_1\circ\cdots\circ t_n$,
\[
\begin{aligned}
\|\theta(y)-t(y)\|&=\|t[\theta(y)-\sum\alpha_it_i(y)]\|\\
&\leqq M\varepsilon/(M+1)<\varepsilon,
\end{aligned}
\]
so $\lim_{t\in T}\|\theta(y)-t(y)\|=0$.\hfill Q.E.D.

In order to apply the theorem to the case of $T=\{t_\pi:\pi\in\mathcal P\}$ acting on $\mathrm{sba}(\unsure{\mathcal A|A},X)$, we must show that for every $\mu$ in $\mathrm{sba}(\unsure{\mathcal A|A},X)$ the set
\[
\{t_\pi\mu:\pi\in\mathcal P\}
\]
is conditionally weakly compact in $\mathrm{sba}(\unsure{\mathcal A|A},X)$. This follows immediately from Proposition 1 below. The theorem then implies that $\mathrm{sba}(\unsure{\mathcal A|A},X)$ can be decomposed into the direct sum of the countably additive members of $\mathrm{sba}(\unsure{\mathcal A|A},X)$ and the purely finitely additive members according to the following definition.
\margin{31}
